\documentclass[aps,preprint]{revtex4}%
\usepackage{amsfonts}
\usepackage{amsmath}
\usepackage{amssymb}
\usepackage[tcidvi]{graphicx}
\usepackage{amsbsy}%
\setcounter{MaxMatrixCols}{30}
%TCIDATA{OutputFilter=latex2.dll}
%TCIDATA{Version=5.50.0.2953}
%TCIDATA{CSTFile=revtex4.cst}
%TCIDATA{Created=Friday, June 20, 2003 15:09:06}
%TCIDATA{LastRevised=Saturday, August 04, 2007 00:20:05}
%TCIDATA{<META NAME="GraphicsSave" CONTENT="32">}
%TCIDATA{<META NAME="SaveForMode" CONTENT="1">}
%TCIDATA{BibliographyScheme=BibTeX}
%TCIDATA{<META NAME="DocumentShell" CONTENT="Articles\SW\REVTeX 4">}
%TCIDATA{Language=American English}
%BeginMSIPreambleData
\providecommand{\U}[1]{\protect\rule{.1in}{.1in}}
%EndMSIPreambleData
\newtheorem{theorem}{Theorem}
\newtheorem{acknowledgement}[theorem]{Acknowledgement}
\newtheorem{algorithm}[theorem]{Algorithm}
\newtheorem{axiom}[theorem]{Axiom}
\newtheorem{claim}[theorem]{Claim}
\newtheorem{conclusion}[theorem]{Conclusion}
\newtheorem{condition}[theorem]{Condition}
\newtheorem{conjecture}[theorem]{Conjecture}
\newtheorem{corollary}[theorem]{Corollary}
\newtheorem{criterion}[theorem]{Criterion}
\newtheorem{definition}[theorem]{Definition}
\newtheorem{example}[theorem]{Example}
\newtheorem{exercise}[theorem]{Exercise}
\newtheorem{lemma}[theorem]{Lemma}
\newtheorem{notation}[theorem]{Notation}
\newtheorem{problem}[theorem]{Problem}
\newtheorem{proposition}[theorem]{Proposition}
\newtheorem{remark}[theorem]{Remark}
\newtheorem{solution}[theorem]{Solution}
\newtheorem{summary}[theorem]{Summary}
\newenvironment{proof}[1][Proof]{\noindent\textbf{#1.} }{\ \rule{0.5em}{0.5em}}
\begin{document}
\preprint{ }
\title[Antisym. Mat]{Canonical Form of Antisymmetric Matrices}
\author{B. Weiner}
\affiliation{Department of Physics, Pennsylvania State University, DuBois PA 15801}
\author{J. V. Ortiz}
\affiliation{Department of Chemistry, Kansas State University, Manhattan, KS 66506-3701}

\pacs{}

\begin{abstract}
The canonical geminal coefficients often advertised as "canonical" for a given
geminal are in fact very non unique. In this paper we show by carefully
examining a Theorem by B. Zumino how to construct unique canonical
coefficients for a geminal. This result leads to expansions of Antisymmetrised
Geminal Power States solely in terms of Natural Generalized Spin Orbitals and
geminal occupation numbers.

\end{abstract}
\date{06/24/2003}
\startpage{1}
\maketitle
\tableofcontents


\section{Introduction\label{intro}}

Antisymmetric Geminal Power (AGP) states provide an unbiased way to treat
correlation for both strongly and weakly correlated systems. This has been
displayed by their use as a reference state for the polarization propagator
\cite{weiner87}\cite{Weiner85}\cite{Weiner84a}\cite{Kurtz83}\cite{Jorgen83} in
the simultaneous production of dissociation curves for several low lying
excited states of light diatomic molecules which were in good agreement with
experimental results. Although AGP states are rarely the exact solutions of
the Schr\"{o}dinger equation they can serve as an unbiased approximation
ansatz generalizing the Independent Particle State (IPS) approach by including
a better description of electron - electron interaction, while at the same
time maintaining a balanced and physically interpretable model of both
molecular and extended systems. In this paper we in fact show that an
$2N$-electron AGP\ state is completely determined by geminal occupation
numbers and Natural General Spin Orbitals (NGSOs), two quantities that have a
clear physical significance. A $2N$-electron AGP\ state is the antisymmetrized
product of $N$ identical $2$-electron geminals thus the properties and form of
the geminal determine the characteristics of the product\ state e.g. the
expectation value expressions and will form the focus of this paper.

The methods used to energy optimize an AGP state in the previously cited
references were not numerically robust nor easily extendable to large systems.
Recently we have \cite{Weiner&Ortiz02} described a method that overcomes these
numerical difficulties and can be applied to molecules of various sizes. A
salient feature of this procedure is that the geminal parameters always refer
to a \emph{fixed} localized atomic orbital basis, which allows one to
visualize the AGP state in a chemical intuitive manner and introduce
approximations to the integral matrix in a natural way. At the same time the
various expressions for the energy and its gradient are more conveniently
formulated in terms of the canonical form (canonical orbitals and
coefficients) of the geminal, which was first discussed in a quantum chemical
context in the works of \cite{Coleman63},\cite{Lowdin56}.

Producing the canonical form of a geminal is equivalent to bringing an
antisymmetric \emph{complex} matrix to a canonical form. This however is not a
straightforward task as such matrices are \emph{not }normal \cite{Encyclo}
(therefore one cannot assume they are diagonalizable) and the methods used for
real antisymmetric matrices have to be carefully modified. However Zumino
\cite{Zumino62} constructively proved a theorem showing how to accomplish this
transformation. In nearly all discussions of the properties of a geminal
various subtleties of the canonical form have been missed or not focused on,
in particular there is \emph{no unique} set of canonical coefficients nor
canonical orbitals. The latter observation is not surprising and since the
canonical orbitals have always known to be indeterminate up to at least
$2\times2$ special unitary transformations, nevertheless the fact that there
is not a unique set of canonical coefficients has not been widely advertised
nor studied in depth. However a corollary in the Zumino paper shows how one
can define a unique set of canonical coefficients as the positive square roots
of the eigenvalues of the First Order Reduced Density Operator (FORDO) of the
$2$-particle geminal state, which in turn defines canonical orbitals up to
special unitary transformations that mix NGSOs associated with a given
occupation number.

After the introduction (section \ref{intro}) we define in section
\ref{geminal} the geminal and reduced density operator associated with this
$2$ -particle state and we also review some aspects of tensor product
expansions that can lead to some ambiguity in the normalization of the geminal
expansion coefficients, in section \ref{zumino} we develop Zuminos theorem and
two corollaries in detail, while in section \ref{numerical} we describe a
procedure for producing the physically significant canonical coefficients and
orbitals and examine the subtleties of numerically implementing this procedure
and we conclude with a discussion in section \ref{discussion}.

\section{Geminals and Antisymmetric Complex Matrices\label{geminal}}

In this section we describe how one forms the matrix of Natural Generalized
Spin Orbitals (NGSO) that brings the complex antisymmetric matrix of geminal
coefficients $\mathbf{c}$\textbf{ }to block diagonal form (see section
\ref{zumino}). It is important to note the significance of the term
"Generalized" which means that the orbitals have both alpha and beta spin
components. This is necessary on two fronts, one mathematical in order to
obtain the correct pairing associated with antisymmetry and one practical in
order to form the most flexible approximation ansatz. In order to produce the
NGSOs one needs to construct the FORDO of the $2$ -particle geminal state
$\left\vert g\right\rangle $ and then sum the indices relating to one of the
particles.\textbf{ }This operation is straightforward in terms of a general
tensor product basis $\left(  \text{see Appendix \ref{tensor}}\right)  $, but
introduces some subtleties of normalization when one uses the standard
expression of the geminal in terms of an antisymmetrized tensor product basis.
For the purposes of clarity we compare the use of these two expansions.

A normalized geminal $\left\vert g\right\rangle $ can be expanded in terms of
a general tensor product orthonormal basis $\left\{  \left\vert \varphi
_{j_{1}}\otimes\varphi_{j_{2}}\right\rangle ;1\leq j_{1},j_{2}\leq r\right\}
$ as
\begin{equation}
\left\vert g\right\rangle =\frac{1}{2}\sum_{1\leq j_{1},j_{2}\leq r}%
c_{j_{1}j_{2}}\left\vert \varphi_{j_{1}}\otimes\varphi_{j_{2}}\right\rangle
;\quad c_{j_{1}j_{2}}=-c_{j_{2}j_{1}},\quad c_{j_{1}j_{2}}\in\mathbb{C}%
\end{equation}
where $\left\{  \left\vert \varphi_{j}\right\rangle ;1\leq j\leq r\right\}  $
is a complete orthonormal basis for the one particle Hilbert space
$\mathcal{H}^{1},$ and $\mathbb{C}$ denotes the complex numbers. The
associated density operator to the state $\left\vert g\right\rangle $ is
\begin{equation}
\left\vert g\right\rangle \left\langle g\right\vert =\frac{1}{4}%
\sum_{\substack{1\leq j_{1},j_{2}\leq r\\1\leq k_{1},k_{2}\leq r}%
}c_{j_{1}j_{2}}c_{k_{1}k_{2}}^{\ast}\left\vert \varphi_{j_{1}}\otimes
\varphi_{j_{2}}\right\rangle \left\langle \varphi_{k_{1}}\otimes\varphi
_{k_{2}}\right\vert
\end{equation}
As $\left\langle g|g\right\rangle =1$ one observes that
\begin{equation}
\operatorname{Tr}\left\{  \mathbf{cc}^{\dagger}\right\}  =4
\end{equation}
The expansion in terms of the tensor product basis $\left\{  \left\vert
\varphi_{j_{1}}\otimes\varphi_{j_{2}}\right\rangle \right\}  $ can be compared
to
\begin{equation}
\left\vert g\right\rangle =\frac{1}{2}\sum_{1\leq j_{1},j_{2}\leq r}\tilde
{c}_{j_{1}j_{2}}\left\vert \varphi_{j_{1}}\wedge\varphi_{j_{2}}\right\rangle
;\quad\tilde{c}_{j_{1}j_{2}}=-\tilde{c}_{j_{2}j_{1}}%
\end{equation}
in terms of the antisymmetrized tensor product basis $\left\{  \left\vert
\varphi_{j_{1}}\wedge\varphi_{j_{2}}\right\rangle ;1\leq j_{1}<j_{2}\leq
r\right\}  $ where
\begin{equation}
\left\vert \varphi_{j_{1}}\wedge\varphi_{j_{2}}\right\rangle =\frac{1}%
{\sqrt{2}}\left(  \left\vert \varphi_{j_{1}}\otimes\varphi_{j_{2}%
}\right\rangle -\left\vert \varphi_{j_{2}}\otimes\varphi_{j_{1}}\right\rangle
\right)
\end{equation}
Expressing the antisymmetric expansion in terms of the general one we obtain
\begin{align}
\frac{1}{2}\sum_{1\leq j_{1},j_{2}\leq r}\tilde{c}_{j_{1}j_{2}}\left\vert
\varphi_{j_{1}}\wedge\varphi_{j_{2}}\right\rangle  &  =\frac{1}{2}\sum_{1\leq
j_{1},j_{2}\leq r}\tilde{c}_{j_{1}j_{2}}\frac{1}{\sqrt{2}}\left(  \left\vert
\varphi_{j_{1}}\otimes\varphi_{j_{2}}\right\rangle -\left\vert \varphi_{j_{2}%
}\otimes\varphi_{j_{1}}\right\rangle \right) \nonumber\\
&  =\frac{1}{2}\frac{1}{\sqrt{2}}\sum_{1\leq j_{1},j_{2}\leq r}\tilde
{c}_{j_{1}j_{2}}\left\vert \varphi_{j_{1}}\otimes\varphi_{j_{2}}\right\rangle
-\frac{1}{2}\frac{1}{\sqrt{2}}\sum_{1\leq j_{1},j_{2}\leq r}\tilde{c}%
_{j_{1}j_{2}}\left\vert \varphi_{j_{2}}\otimes\varphi_{j_{1}}\right\rangle
\nonumber\\
&  =\frac{1}{2}\frac{1}{\sqrt{2}}\sum_{1\leq j_{1},j_{2}\leq r}\tilde
{c}_{j_{1}j_{2}}\left\vert \varphi_{j_{1}}\otimes\varphi_{j_{2}}\right\rangle
+\frac{1}{2}\frac{1}{\sqrt{2}}\sum_{1\leq j_{1},j_{2}\leq r}\tilde{c}%
_{j_{1}j_{2}}\left\vert \varphi_{j_{1}}\otimes\varphi_{j_{2}}\right\rangle
\nonumber\\
&  =\frac{2}{\sqrt{2}}\sum_{1\leq j_{1},j_{2}\leq r}\tilde{c}_{j_{1}j_{2}%
}\left\vert \varphi_{j_{1}}\otimes\varphi_{j_{2}}\right\rangle
\end{align}
which shows that
\begin{equation}
c_{j_{1}j_{2}}=\sqrt{2}\tilde{c}_{j_{1}j_{2}}%
\end{equation}


The contraction map \textrm{L}$_{2}^{1}$ \cite{Weiner84}\cite{kummar}%
\cite{lowdin} acting on $\left\vert g\right\rangle \left\langle g\right\vert $
produces
\begin{align}
\mathrm{L}_{2}^{1}\left(  \left\vert g\right\rangle \left\langle g\right\vert
\right)   &  =\frac{1}{4}\sum_{\substack{1\leq j_{1},j_{2}\leq r\\1\leq
k_{1},k_{2}\leq r}}c_{j_{1}j_{2}}c_{k_{1}k_{2}}^{\ast}\left\langle
\varphi_{k_{2}}|\varphi_{j_{2}}\right\rangle \left\vert \varphi_{j_{1}%
}\right\rangle \left\langle \varphi_{k_{1}}\right\vert =\frac{1}{4}%
\sum_{\substack{1\leq j_{1},j_{2}\leq r\\1\leq k_{1},k_{2}\leq r}%
}c_{j_{1}j_{2}}c_{k_{1}k_{2}}^{\ast}\delta_{k_{2}j_{2}}\left\vert
\varphi_{j_{1}}\right\rangle \left\langle \varphi_{k_{1}}\right\vert
\nonumber\\
&  =\frac{1}{4}\sum_{\substack{1\leq j_{1},j_{2}\leq r\\1\leq k_{1}\leq
r}}c_{j_{1}j_{2}}c_{k_{1}j_{2}}^{\ast}\left\vert \varphi_{j_{1}}\right\rangle
\left\langle \varphi_{k_{1}}\right\vert =\frac{1}{4}\sum_{1\leq j_{1}%
,j_{2}\leq r}\left(  \mathbf{cc}^{\dagger}\right)  _{j_{1}k_{1}}\left\vert
\varphi_{j_{1}}\right\rangle \left\langle \varphi_{k_{1}}\right\vert
\end{align}
The First Order Reduced Density Operator (FORDO) $D^{1}\left(  g\right)  $
normalized to the number of electrons in the system, (in this case $2)$, is
then given by
\begin{equation}
D^{1}\left(  g\right)  =2\mathrm{L}_{2}^{1}\left(  \left\vert g\right\rangle
\left\langle g\right\vert \right)
\end{equation}
and the First Order Reduced Density Matrix (FORDM), $\mathbb{D}_{\varphi}%
^{1}\left(  g\right)  ,$ which is the \emph{representation} of $D^{1}\left(
g\right)  $ in the $\left\{  \mathbf{\varphi}\right\}  $ basis by
\begin{equation}
\mathbb{D}_{\varphi}^{1}\left(  g\right)  =\frac{1}{2}\mathbf{cc}^{\dagger
}=\mathbf{\tilde{c}\tilde{c}}^{\dagger}%
\end{equation}
both of which have trace equal to $2$ i.e.
\begin{equation}
\operatorname{Tr}\left\{  D^{1}\left(  g\right)  \right\}  =\frac{1}%
{2}\operatorname{Tr}\left\{  \mathbf{cc}^{\dagger}\right\}  =\operatorname{Tr}%
\left\{  \mathbf{\tilde{c}\tilde{c}}^{\dagger}\right\}  =2
\end{equation}
One can also observe that
\begin{equation}
\mathbb{D}_{\varphi}^{1}\left(  g\right)  ^{\ast}=\frac{1}{2}\mathbf{c}%
^{\dagger}\mathbf{c}%
\end{equation}
and as $\mathbf{c}$ is antisymmetric
\begin{equation}
\mathbb{D}_{\varphi}^{1}\left(  g\right)  =\frac{1}{2}\mathbf{cc}^{\dagger
}=-\frac{1}{2}\mathbf{cc}^{\ast}=\frac{1}{2}\mathbf{c}^{t}\mathbf{c}^{\ast}%
\end{equation}
The NGSO's $\left\{  \left\vert y_{j}\right\rangle ;1\leq j\leq r\right\}  $
of $D^{1}\left(  g\right)  $ have expansion coefficients $\left\{
\mathbf{y}_{j};1\leq j\leq r\right\}  $ in the $\left\{  \mathbf{\varphi
}\right\}  $ basis that are the eigenvectors of $\mathbb{D}_{\varphi}%
^{1}\left(  g\right)  $ given by
\begin{equation}
\mathbb{D}_{\varphi}^{1}\left(  g\right)  \mathbf{y}_{j}=n_{j}\mathbf{y}_{j}%
\end{equation}
leading to
\begin{equation}
\mathbb{Y}^{\dagger}\mathbb{D}_{\varphi}^{1}\left(  g\right)  \mathbb{Y=}%
\left(
\begin{array}
[c]{cc}%
\mathbf{n} & \mathbf{0}\\
\mathbf{0} & \mathbf{n}%
\end{array}
\right)
\end{equation}
where
\begin{equation}
\mathbf{n=}\left(
\begin{array}
[c]{ccc}%
n_{1} & 0 & 0\\
0 & \ddots & 0\\
0 & 0 & n_{s}%
\end{array}
\right)  ;\qquad n_{j}\geq0,\quad r=2s
\end{equation}
and
\begin{equation}
\mathbb{Y}^{\dagger}\mathbb{Y}=\mathbb{YY}^{\dagger}=\mathbb{I}_{r}%
\end{equation}
Although this transformation brings $\mathbb{D}_{\varphi}^{1}\left(  g\right)
$ to diagonal form it does $\emph{not}$ in general bring $\mathbf{c}$ to a
canonical form. In the next section we show, how using a Theorem due to Zumino
\cite{Zumino62} one can achieve this result.

\section{Zuminos' Theorem\label{zumino}}

In this section we first describe Zuminos' Theorem and its proof in detail and
then discuss an important corollary that leads to the definition of a unique,
physically meaningful, canonical form of $\mathbf{c.}$

\begin{theorem}
If $\mathbf{c}$ is a $2s\times2s$ complex antisymmetric matrix then there
exists a unitary matrix $\mathbb{T}$ such that
\begin{equation}
\mathbb{T}\mathbf{c}\mathbb{T}^{t}\mathbb{=}\left(
\begin{array}
[c]{cc}%
\mathbf{0} & \mathbf{\zeta}\\
-\mathbf{\zeta} & \mathbf{0}%
\end{array}
\right)
\end{equation}
where
\begin{equation}
\mathbf{\zeta=}\left(
\begin{array}
[c]{ccc}%
\zeta_{1} & 0 & 0\\
0 & \ddots & 0\\
0 & 0 & \zeta_{s}%
\end{array}
\right)  ;\qquad\zeta_{j}\in\mathbb{C},\quad1\leq j\leq s
\end{equation}


\begin{proof}
We first note some relationships implied by the antisymmetry of $\mathbf{c}$%
\begin{equation}
\mathbf{c}^{\dagger}\mathbf{c\,=-c}^{\ast}\mathbf{c}=\left(  \mathbf{cc}%
^{\dagger}\right)  ^{t}=\mathbf{\,}\left(  \mathbf{cc}^{\dagger}\right)
^{\ast} \label{cc}%
\end{equation}
(so $\mathbf{c}$ is not a normal operator in general) and one by the
associativity of matrix product
\begin{equation}
\mathbf{c}\left(  \mathbf{c}^{\dagger}\mathbf{c}\right)  =\left(
\mathbf{cc}^{\dagger}\right)  \mathbf{c}%
\end{equation}
which by using eq. $\left(  \ref{cc}\right)  $ gives
\begin{equation}
\left(  \mathbf{cc}^{\dagger}\right)  \mathbf{c=c}\left(  \mathbf{cc}%
^{\dagger}\right)  ^{\ast} \label{cr}%
\end{equation}
Then we introduce $\mathbf{c}_{U}\mathbf{c}_{U}^{\dagger}$ through the
definition
\begin{align}
\mathbf{c}  &  =\mathbb{U}\mathbf{c}_{U}\mathbb{U}^{t}\\
\mathbf{c}_{U}  &  =\mathbb{U}^{\dagger}\mathbf{c}\mathbb{U}^{\ast}%
\end{align}
where $\mathbb{U}$ is \emph{any} unitary matrix and obtain that
\begin{align}
\mathbf{cc}^{\dagger}  &  =\mathbb{U}\mathbf{c}_{U}\mathbb{U}^{t}%
\mathbb{U}^{\ast}\mathbf{c}_{U}^{\dagger}\mathbb{U}^{\dagger}=\mathbb{U}%
\mathbf{c}_{U}\mathbf{c}_{U}^{\dagger}\mathbb{U}^{\dagger}\nonumber\\
\mathbf{c}_{U}\mathbf{c}_{U}^{\dagger}  &  =\mathbb{U}^{\dagger}%
\mathbf{cc}^{\dagger}\mathbb{U}%
\end{align}
and as $\mathbf{c}_{U}$ is also antisymmetric
\begin{equation}
\mathbf{c}_{U}\mathbf{c}_{U}^{\dagger}=-\mathbf{c}_{U}\mathbf{c}_{U}^{\ast
}=\left(  \mathbf{c}_{U}^{\dagger}\mathbf{c}_{U}\right)  ^{t}=\mathbf{\,}%
\left(  \mathbf{c}_{U}^{\dagger}\mathbf{c}_{U}\right)  ^{\ast}.
\end{equation}
Hence we have that
\begin{equation}
\left(  \mathbf{cc}^{\dagger}\right)  \mathbf{c-c}\left(  \mathbf{cc}%
^{\dagger}\right)  ^{\ast}=\left(  \mathbf{c}_{U}\mathbf{c}_{U}^{\dagger
}\right)  \mathbf{c}_{U}-\mathbf{c}_{U}\left(  \mathbf{c}_{U}\mathbf{c}%
_{U}^{\dagger}\right)  ^{\ast}=0 \label{cu}%
\end{equation}
We can choose $\mathbb{U}$ to diagonalize the hermitian matrix $\mathbf{cc}%
^{\dagger}$ i.e.%
\begin{equation}
\mathbb{U}^{\dagger}\mathbf{cc}\mathbb{^{\dagger}U}=\left(
\begin{array}
[c]{ccc}%
n_{1} & 0 & 0\\
0 & \ddots & 0\\
0 & 0 & n_{r}%
\end{array}
\right)  =\mathbf{c}_{U}\mathbf{c}_{U}^{\dagger};\qquad n_{r}\geq n_{r-1}%
\geq\cdots\geq n_{1}\geq0.
\end{equation}
Now using the relationships in eq. $\left(  \ref{cu}\right)  $ for this
$\mathbf{c}_{U}$ we obtain%
\begin{equation}
\sum_{1\leq j\leq r}\left\{  n_{i}\delta_{ij}c_{Ujk}-c_{Uij}n_{j}\delta
_{jk}\right\}  =0\Rightarrow\left(  n_{i}-n_{k}\right)  c_{Uik}=0;\qquad1\leq
i,k\leq r \label{dub1}%
\end{equation}
which shows that
\begin{equation}
n_{k}\neq n_{i}\Rightarrow c_{Uik}=0;\qquad\qquad1\leq i\neq k\leq r
\end{equation}
The diagonal elements
\begin{equation}
c_{Uii}=0;\qquad\qquad1\leq i\leq r
\end{equation}
are always zero as $\mathbf{c}_{U}$ is antisymmetric. Thus the matrix
$\mathbf{c}_{U}$ breaks into submatrices $\left\{  \mathbb{C}_{\kappa}%
;1\leq\kappa\leq p\right\}  $ corresponding to groups of $n^{\prime}s$ that
are equal. Consider the partitioning of $\left\{  1,\ldots,r\right\}  $ into
$p$ groups that correspond to equal $n^{\prime}s$ i.e. $\left\{  \left\{
1,\ldots,m_{1}\right\}  ,\cdots,\left\{  m_{p-1}+1,\ldots r\right\}  \right\}
$ where
\begin{align}
n_{1}  &  =n_{2}=\cdots=n_{m_{1}}=\eta_{1}\nonumber\\
&  \vdots\nonumber\\
n_{m_{p-1}+1}  &  =n_{m_{p-1}+2}=\cdots=n_{r}=\eta_{p}%
\end{align}
then the matrix $\mathbf{c}_{U}$ has the block diagonal form
\begin{equation}
\mathbf{c}_{U}=\left(
\begin{array}
[c]{cccc}%
\mathbf{c}_{1} & \mathbb{O} & \cdots & \mathbb{O}\\
\mathbb{O} & \ddots & \mathbb{O} & \vdots\\
\vdots & \mathbb{O} & \ddots & \mathbb{O}\\
\mathbb{O} & \cdots & \mathbb{O} & \mathbf{c}_{p}%
\end{array}
\right)
\end{equation}
and
\begin{equation}
\mathbf{c_{\kappa}^{\dagger}c_{\kappa}}=\eta_{\kappa}\mathbb{I}_{d_{\kappa}%
};\quad1\leq\kappa\leq p \label{dub2}%
\end{equation}
where $d_{\kappa}$ is the dimension of $\kappa^{th}$-block.

If $\eta_{\kappa}=0$ then
\begin{equation}
\sum_{k\in\left\{  m_{\kappa-1}+1,\ldots m_{\kappa}\right\}  }c_{\kappa
ik}c_{\kappa jk}^{\ast}=0\qquad\forall i,j\in\left\{  m_{\kappa-1}%
+1,\ldots,m_{\kappa}\right\}
\end{equation}
so in particular
\begin{equation}
\sum_{k\in\left\{  m_{\kappa-1}+1,\ldots m_{\kappa}\right\}  }\left\vert
c_{\kappa jk}\right\vert ^{2}=0\qquad\forall j\in\left\{  m_{\kappa
-1}+1,\ldots,m_{\kappa}\right\}
\end{equation}
which implies that $c_{\kappa jk}=0\quad\forall j,k\in\left\{  m_{\kappa
-1}+1,\ldots,m_{\kappa}\right\}  $.

If $\eta_{\kappa}\neq0$ then $\mathbb{B}_{\kappa}=\frac{1}{\sqrt{\eta_{\kappa
}}}\mathbf{c}_{\kappa}$ is a $d_{\kappa}\times d_{\kappa}$ antisymmetric
unitary matrix
\begin{align}
\mathbb{B}_{\kappa}^{\dagger}\mathbb{B}_{\kappa}  &  =\mathbb{B}_{\kappa
}\mathbb{B}_{\kappa}^{\dagger}=\mathbb{I}_{d_{\kappa}}\nonumber\\
&  \Rightarrow\mathbb{B}_{\kappa}^{\ast}\mathbb{B}_{\kappa}=\mathbb{B}%
_{\kappa}\mathbb{B}_{\kappa}^{\ast}%
\end{align}
i.e. $\left[  \mathbb{B}_{\kappa},\mathbb{B}_{\kappa}^{\dagger}\right]
=\left[  \mathbb{B}_{\kappa},\mathbb{B}_{\kappa}^{\ast}\right]  =0$ showing
that $\mathbb{B}_{\kappa}$ is a normal matrix. One can then define the
matrices
\begin{align}
\mathbb{N}_{\kappa}  &  =\frac{1}{2}\left\{  \mathbb{B}_{\kappa}%
+\mathbb{B}_{\kappa}^{\ast}\right\} \nonumber\\
\mathbb{M}_{\kappa}  &  =\frac{-i}{2}\left\{  \mathbb{B}_{\kappa}%
-\mathbb{B}_{\kappa}^{\ast}\right\}
\end{align}
which are not in general self adjoint but are real antisymmetric and
\emph{commute }with each other so that
\begin{equation}
\mathbb{B}_{\kappa}=\mathbb{N}_{\kappa}+i\mathbb{M}_{\kappa}.
\end{equation}
These matrices can be simultaneously transformed into the canonical forms
\begin{equation}
\left(
\begin{array}
[c]{cc}%
\mathbf{0} & \mathbf{\alpha}_{\kappa}\\
-\mathbf{\alpha}_{\kappa} & \mathbf{0}%
\end{array}
\right)  \text{ and }\left(
\begin{array}
[c]{cc}%
\mathbf{0} & \mathbf{\beta}_{\kappa}\\
-\mathbf{\beta}_{\kappa} & \mathbf{0}%
\end{array}
\right)  \label{alphabeta}%
\end{equation}
by a real orthogonal transformation $\mathbb{O}_{\kappa}$, where
\begin{equation}
\mathbf{\alpha}_{\kappa}=\left(
\begin{array}
[c]{ccc}%
\alpha_{\kappa1} & 0 & 0\\
0 & \ddots & 0\\
0 & 0 & \alpha_{\kappa d_{\kappa}}%
\end{array}
\right)  \text{ and\quad}\mathbf{\beta}_{\kappa}=\left(
\begin{array}
[c]{ccc}%
\beta_{\kappa1} & 0 & 0\\
0 & \ddots & 0\\
0 & 0 & \beta_{\kappa d_{\kappa}}%
\end{array}
\right)  \label{alphabeta1}%
\end{equation}
as shown in Appendix \ref{realanti}.

The same transformation thus acts on $\mathbb{B}_{\kappa}$ to produce
\begin{equation}
\mathbb{O}_{\kappa}^{t}\mathbb{B}_{\kappa}\mathbb{O}_{\kappa}=\mathbb{O}%
_{\kappa}^{t}\left\{  \mathbb{N}_{\kappa}+i\mathbb{M}_{\kappa}\right\}
\mathbb{O}_{\kappa}=\left(
\begin{array}
[c]{cc}%
\mathbf{0} & \mathbf{\alpha}_{\kappa}+i\mathbf{\beta}_{\kappa}\\
-\left(  \mathbf{\alpha}_{\kappa}+i\mathbf{\beta}_{\kappa}\right)  &
\mathbf{0}%
\end{array}
\right)  =\left(
\begin{array}
[c]{cc}%
\mathbf{0} & e^{i\mathbf{\theta}_{\kappa}}\\
-e^{i\mathbf{\theta}_{\kappa}} & \mathbf{0}%
\end{array}
\right)
\end{equation}
and thus%
\begin{equation}
\mathbb{O}_{\kappa}^{t}\mathbb{C}_{\kappa}\mathbb{O}_{\kappa}=\eta_{\kappa
}^{\frac{1}{2}}\left(
\begin{array}
[c]{cc}%
\mathbf{0} & e^{i\mathbf{\theta}_{\kappa}}\\
-e^{i\mathbf{\theta}_{\kappa}} & \mathbf{0}%
\end{array}
\right)  =\left(
\begin{array}
[c]{cc}%
\mathbf{0} & \eta_{\kappa}^{\frac{1}{2}}e^{i\mathbf{\theta}_{\kappa}}\\
-\eta_{\kappa}^{\frac{1}{2}}e^{i\mathbf{\theta}_{\kappa}} & \mathbf{0}%
\end{array}
\right)  =\left(
\begin{array}
[c]{cc}%
\mathbf{0} & \mathbf{\zeta}_{\kappa}\\
-\mathbf{\zeta}_{\kappa} & \mathbf{0}%
\end{array}
\right)
\end{equation}
(which is valid even when $\eta_{\kappa}=0$) $\lceil$Note that $\mathbf{\theta
}_{\kappa}\equiv\left(  \theta_{\kappa1},\ldots,\theta_{\kappa d_{\kappa}%
}\right)  $ and $\theta_{\kappa i}\neq\theta_{\kappa j}$ in general$\rfloor$
where
\begin{equation}
\mathbf{\zeta}_{\kappa}=\left(
\begin{array}
[c]{ccc}%
\zeta_{\kappa1} & 0 & 0\\
0 & \ddots & 0\\
0 & 0 & \zeta_{\kappa d_{\kappa}}%
\end{array}
\right)  .
\end{equation}
This allows us to define
\begin{equation}
\mathbb{O=}\left(
\begin{array}
[c]{ccc}%
\mathbb{O}_{1} & \mathbf{0} & \mathbf{0}\\
\mathbf{0} & \ddots & \mathbf{0}\\
\mathbf{0} & \mathbf{0} & \mathbb{O}_{p}%
\end{array}
\right)
\end{equation}
and see that
\begin{equation}
\mathbb{O}^{t}\mathbf{c}_{U}\mathbb{O}=\left(
\begin{array}
[c]{ccc}%
\begin{array}
[c]{cc}%
\mathbf{0} & \mathbf{\zeta}_{1}\\
-\mathbf{\zeta}_{1} & \mathbf{0}%
\end{array}
& \mathbf{0} & \mathbf{0}\\
\mathbf{0} & \ddots & \mathbf{0}\\
\mathbf{0} & \mathbf{0} &
\begin{array}
[c]{cc}%
\mathbf{0} & \mathbf{\zeta}_{p}\\
-\mathbf{\zeta}_{p} & \mathbf{0}%
\end{array}
\end{array}
\right)  \label{cor1}%
\end{equation}
by rearranging the columns of $\mathbb{O}$ to produce $\widetilde{\mathbb{O}}$
we obtain
\begin{equation}
\widetilde{\mathbb{O}}^{t}\mathbf{c}_{U}\widetilde{\mathbb{O}}=\left(
\begin{array}
[c]{cc}%
\mathbf{0} & \mathbf{\zeta}\\
-\mathbf{\zeta} & \mathbf{0}%
\end{array}
\right)
\end{equation}
and hence using $\mathbf{c}_{U}=\mathbb{U}^{\dagger}\mathbf{c}\mathbb{U}%
^{\ast}$ we conclude that
\begin{equation}
\mathbb{T}\mathbf{c}\mathbb{T}^{t}=\widetilde{\mathbb{O}}^{t}\mathbb{U}%
^{\dagger}\mathbf{c}\mathbb{U}^{\ast}\widetilde{\mathbb{O}}=\left(
\begin{array}
[c]{cc}%
\mathbf{0} & \mathbf{\zeta}\\
-\mathbf{\zeta} & \mathbf{0}%
\end{array}
\right)
\end{equation}
and define $\mathbb{T}$ as the unitary matrix
\begin{equation}
\mathbb{T}=\widetilde{\mathbb{O}}^{t}\mathbb{U}^{\dagger}%
\end{equation}
where $\mathbb{TT}^{\dagger}=\mathbb{T^{\dagger}T=I}_{r}$.
\end{proof}
\end{theorem}

We now derive an important corollary that will be very useful in the physical
interpretation of the density operators and expectation value expressions of AGP\ states:

\begin{corollary}
If $\mathbf{c}$ is a $2s\times2s$ complex antisymmetric matrix then there
exists a unitary matrix $\mathbb{W}$ such that
\begin{equation}
\mathbb{W}\mathbf{c}\mathbb{W}^{t}\mathbb{=}\left(
\begin{array}
[c]{cc}%
\mathbf{0} & \mathbf{n}^{\frac{1}{2}}\\
-\mathbf{n}^{\frac{1}{2}} & \mathbf{0}%
\end{array}
\right)
\end{equation}
where
\begin{equation}
\mathbf{n}^{\frac{1}{2}}=\left(
\begin{array}
[c]{ccc}%
n_{1}^{\frac{1}{2}} & 0 & 0\\
0 & \ddots & 0\\
0 & 0 & n_{s}^{\frac{1}{2}}%
\end{array}
\right)  ;\qquad n_{j}^{\frac{1}{2}}\in\mathbb{R},\quad1\leq j\leq s
\end{equation}
and $\left\{  n_{j}\right\}  $ are the eigenvalues of $\mathbf{cc}^{\dagger}$.
\end{corollary}

\begin{proof}
Using eq. $\left(  \ref{cor1}\right)  $ we can define $\mathbb{W}$ as
\begin{equation}
\mathbb{W=}\left(
\begin{array}
[c]{cccccc}%
e^{-i\frac{\mathbf{\theta}_{1}}{2}} & \mathbf{0} & \cdots & \cdots &
\mathbf{0} & \mathbf{0}\\
\mathbf{0} & e^{-i\frac{\mathbf{\theta}_{1}}{2}} & \mathbf{0} & \cdots &
\mathbf{0} & \mathbf{0}\\
\mathbf{\vdots} & \mathbf{0} & \mathbf{\ddots} & \cdots & \mathbf{\vdots} &
\mathbf{\vdots}\\
\mathbf{\vdots} & \mathbf{\vdots} & \mathbf{0} & \ddots & \mathbf{0} &
\mathbf{\vdots}\\
\mathbf{0} & \mathbf{0} & \cdots & \mathbf{0} & e^{-i\frac{\mathbf{\theta}%
_{p}}{2}} & \mathbf{0}\\
\mathbf{0} & \mathbf{0} & \cdots & \cdots & \mathbf{0} & e^{-i\frac
{\mathbf{\theta}_{p}}{2}}%
\end{array}
\right)  \mathbb{T}%
\end{equation}
where
\begin{equation}
\mathbb{W}^{\dagger}\mathbb{W=WW}^{\dagger}\mathbb{=I}_{2s}%
\end{equation}
and
\begin{equation}
\mathbf{c}=\mathbb{W}^{\dagger}\left(
\begin{array}
[c]{cc}%
\mathbf{0} & \mathbf{n}^{\frac{1}{2}}\\
-\mathbf{n}^{\frac{1}{2}} & \mathbf{0}%
\end{array}
\right)  \mathbb{W}^{\ast}=\mathbb{T}^{\dagger}\left(
\begin{array}
[c]{cc}%
\mathbf{0} & \mathbf{\zeta}\\
-\mathbf{\zeta} & \mathbf{0}%
\end{array}
\right)  \mathbb{T}^{\ast} \label{cex}%
\end{equation}

\end{proof}

Eq. $\left(  \ref{cex}\right)  $ shows that the canonical NGSOs $\left\{
\left\vert \psi_{k}\right\rangle \right\}  $ are represented by the columns of
the matrix $\mathbb{W}^{\dagger}$ in the basis $\left\{  \left\vert
\varphi_{k}\right\rangle \right\}  $ giving
\begin{align}
W\left(  \left\vert \mathbf{\psi}\right\rangle \right)   &  =\left(
\left\vert \mathbf{\varphi}\right\rangle \right) \nonumber\\
W^{\dagger}\left(  \left\vert \mathbf{\varphi}\right\rangle \right)   &
=\left(  \left\vert \mathbf{\varphi}\right\rangle \right)  \mathbb{W}%
^{\dagger}=\left(  \left\vert \mathbf{\psi}\right\rangle \right) \nonumber\\
W^{\dagger}\left(  \left\vert \mathbf{\varphi}\right\rangle \right)   &
=\left(  \left\vert \mathbf{\varphi}\right\rangle \right)  \left(
\begin{array}
[c]{ccc}%
\mathbf{w}_{1}^{\dagger} & \cdots & \mathbf{w}_{2s}^{\dagger}%
\end{array}
\right)
\end{align}
where
\begin{align}
\left\vert g\right\rangle  &  =\frac{1}{2}\sum_{1\leq k_{1},k_{2}\leq
r}c_{k_{1}k_{2}}\left\vert \varphi_{k_{1}}\varphi_{k_{2}}\right\rangle
=\frac{1}{2}\sum_{1\leq k_{1},k_{2}\leq r}\left(  \mathbb{W}^{\dagger}\left(
\begin{array}
[c]{cc}%
\mathbf{0} & \mathbf{n}^{\frac{1}{2}}\\
-\mathbf{n}^{\frac{1}{2}} & \mathbf{0}%
\end{array}
\right)  \mathbb{W}^{\ast}\right)  _{k_{1}k_{2}}\left\vert \varphi_{k_{1}%
}\varphi_{k_{2}}\right\rangle \nonumber\\
&  =\sum_{1\leq k_{1}\leq s}n_{k_{1}}^{\frac{1}{2}}\left\vert W^{\dagger
}\varphi_{k_{1}}W^{\dagger}\varphi_{k_{1}+s}\right\rangle =\sum_{1\leq
k_{1}\leq s}n_{k_{1}}^{\frac{1}{2}}\left\vert \psi_{k_{1}}\psi_{k_{1}%
+s}\right\rangle
\end{align}
One can obtain a non unique canonical expansion by setting
\begin{align}
&  \left\vert g\right\rangle =\sum_{1\leq k_{1}\leq s}n_{k_{1}}^{\frac{1}{2}%
}\left\vert e^{i\frac{\mathbf{\theta}_{k_{1}}}{2}}T^{\dagger}\varphi_{k_{1}%
}e^{i\frac{\mathbf{\theta}_{k_{1}}}{2}}T^{\dagger}\varphi_{k_{1}%
+s}\right\rangle =\sum_{1\leq k_{1}\leq s}n_{k_{1}}^{\frac{1}{2}}%
e^{i\theta_{k_{1}}}\left\vert T^{\dagger}\varphi_{k_{1}}T^{\dagger}%
\varphi_{k_{1}+s}\right\rangle \nonumber\\
&  \left.  \qquad\qquad\qquad\qquad\qquad\qquad\qquad\qquad\qquad\right.
=\sum_{1\leq k_{1}\leq s}n_{k_{1}}^{\frac{1}{2}}e^{i\theta_{k_{1}}}\left\vert
\chi_{k_{1}}\chi_{k_{1}+s}\right\rangle
\end{align}
so that
\begin{align}
T\left(  \left\vert \mathbf{\chi}\right\rangle \right)   &  =\left(
\left\vert \mathbf{\varphi}\right\rangle \right) \nonumber\\
T^{\dagger}\left(  \left\vert \mathbf{\chi}\right\rangle \right)   &  =\left(
\left\vert \mathbf{\varphi}\right\rangle \right)  \mathbb{T}^{\dagger}=\left(
\left\vert \mathbf{\chi}\right\rangle \right)
\end{align}
The non uniqueness of the complex canonical coefficients is easily seen by
noting that
\begin{equation}
\left\vert g\right\rangle =\sum_{1\leq k_{1}\leq s}n_{k_{1}}^{\frac{1}{2}%
}e^{i\theta_{k_{1}}}e^{-i\nu_{k_{1}}}\left\vert e^{i\frac{\nu_{k_{1}}}{2}}%
\chi_{k_{1}}e^{i\frac{\nu_{k_{1}}}{2}}\chi_{k_{1}+s}\right\rangle =\sum_{1\leq
k_{1}\leq s}n_{k_{1}}^{\frac{1}{2}}e^{i\left(  \theta_{k_{1}}-\nu_{k_{1}%
}\right)  }\left\vert \tilde{\chi}_{k_{1}}\tilde{\chi}_{k_{1}+s}\right\rangle
\end{equation}
where
\begin{equation}
\left\vert \tilde{\chi}_{k_{1}}\right\rangle =\left\vert e^{i\frac{\nu_{k_{1}%
}}{2}}\chi_{k_{1}}\right\rangle
\end{equation}


We now prove the often stated fact that the eigenvalues of $\mathbb{D}%
_{\varphi}^{1}\left(  g\right)  $ are at least doubly degenerate

\begin{corollary}
The matrix $\mathbf{cc}^{\dagger}$ has eigenvalues that are degenerate of
degree $2m;$ $1\leq m\leq s$.

\begin{proof}
Eq. $\left(  \ref{dub1}\right)  $ shows that if the eigenvalue $n_{j}$ is non
degenerate then $c_{Ujk}=0;1\leq k\leq r$. Eq. $\left(  \ref{dub2}\right)  $
shows in particular that for such non degenerate eigenvalues $\left\vert
c_{Ujj}\right\vert ^{2}=n_{j}=0$, thus all non degenerate eigenvalues must be
equal to $0.$

If an eigenvalue is doubly degenerate then the associated unitary
antisymmetric block $\mathbb{B}$ is $2\times2$ and it is unitary and
antisymmetric, requirements that it can easily be satisfied by non zero
$\mathbb{B}$'s. This is true for all evenly degenerate eigenvalues.

If an eigenvalue is triply degenerate then the associated unitary
antisymmetric block $\mathbb{B}$ is $3\times3$. As it is antisymmetric one or
all three of its eigenvalues must be $0$ (as by Appendix \ref{pairing} the
eigenvalues must be +/- pairs or zero) which is a contradiction as the matrix
$\mathbb{B}$ is also unitary so cannot have zero eigenvalues. Thus the block
cannot be $3\times3,$ the same argument holds for any odd dimensioned block
hence only even dimensional blocks can exist.
\end{proof}
\end{corollary}

\section{Numerical Construction of Canonical GSO\label{numerical}}

The first step is to construct the FORDM $\mathbb{D}_{\varphi}^{1}\left(
g\right)  =\frac{1}{2}\mathbf{cc}^{\dagger}$ then diagonalize it by a unitary
transformation $\mathbb{Y}$. One then transforms $\mathbf{c}$ by
$\mathbb{Y}^{\dagger}\mathbf{c}\mathbb{Y}^{\ast}$ into block diagonal form
with blocks $\left\{  \mathbb{C}_{\kappa}\right\}  $ each associated with a
degenerate eigenvalue $\eta_{\kappa}$ of $\mathbb{D}_{\varphi}^{1}\left(
g\right)  .$ A unitary antisymmetric matrix $\mathbb{B}_{\kappa}$ is then
produced from each block $\mathbb{C}_{\kappa}$ associated with a non zero
eigenvalue by
\begin{equation}
\mathbb{B}_{\kappa}=\frac{1}{\sqrt{\eta_{\kappa}}}\mathbf{c}_{\kappa}%
\end{equation}
The next step is to produce a unitary transformation that diagonalizes
$\mathbb{B}_{\kappa}$. One procedure to achieve this would be simultaneously
diagonalize the self adjoint real and imaginary parts of $\mathbb{B}$
\begin{align}
\mathbb{B}_{R}  &  =\operatorname{Re}\mathbb{B}=\frac{1}{2}\left\{
\mathbb{B+B}^{\dagger}\right\} \nonumber\\
\mathbb{B}_{I}  &  =\operatorname{Im}\mathbb{B}=-\frac{i}{2}\left\{
\mathbb{B-B}^{\dagger}\right\}
\end{align}
which are both antisymmetric
\begin{equation}
\mathbb{B}_{R}=-\mathbb{B}_{R}^{t}\text{ and }\mathbb{B}_{I}=-\mathbb{B}%
_{I}^{t}%
\end{equation}
and commute with each other
\begin{equation}
\left[  \mathbb{B}_{R},\mathbb{B}_{I}\right]  =0
\end{equation}
where we have dropped the subscript $\kappa$. As neither matrix is positive
definite the standard way to do this is to first diagonalize one, then
diagonalize the other in the eigen basis of the first. However if the
dimension of the nonzero blocks $\left\{  \mathbb{B}_{R},\mathbb{B}%
_{I}\right\}  $ are greater than $2$ e.g. $4,6,$ $\cdots,$ etc., there are
more than $1$ positive and $1$ negative eigenvalues (there are always equal
numbers of positive and negative eigenvalues as $\left\{  \mathbb{B}%
_{R},\mathbb{B}_{I}\right\}  $ are antisymmetric), and numerical
diagonalization routines will order these positive and negative sets in some
predetermined sequence. We know from eqs. $\left(  \ref{alphabeta}\right)  $
and $\left(  \ref{alphabeta1}\right)  $ and Appendix \ref{realanti} that in
order to obtain a single unitary transformation that produces a single
orthogonal transformation $\mathbb{O}$ that simultaneously brings $\left\{
\mathbb{B}_{R},\mathbb{B}_{I}\right\}  $ to a canonical form the eigenvalues
$\left\{  \mp i\alpha_{j},\mp i\beta_{j};1\leq j\leq d\right\}  $ of
$\mathbb{B}_{R}$ and $\mathbb{B}_{I}$ must correspond in a very precise
manner, so one would need to find and apply some principle of correspondence
and a procedure that to accomplish this.

One can avoid this problem by directly diagonalizing the antisymmetric normal,
(therefore diagonalizable), unitary matrix $\mathbb{B}$ with a unitary matrix
$\mathbb{V}$ to obtain
\begin{equation}
\mathbb{B}=\left(  \mathbf{v}_{1}\cdots\mathbf{v}_{d}\mathbf{v}_{1}^{\ast
}\mathbf{\cdots v}_{d}^{\ast}\right)  \left(
\begin{array}
[c]{cccccc}%
-ie^{i\theta_{1}} & 0 & 0 & 0 & \cdots & 0\\
0 & \ddots & 0 & 0 & \cdots & 0\\
0 & \cdots & -ie^{i\theta_{d}} & 0 & \cdots & 0\\
0 & \cdots & 0 & ie^{i\theta_{d+1}} & \cdots & 0\\
0 & \cdots & 0 & 0 & \ddots & 0\\
0 & \cdots & 0 & 0 & \cdots & ie^{i\theta_{2d}}%
\end{array}
\right)  \left(
\begin{array}
[c]{c}%
\mathbf{v}_{1}\\
\vdots\\
\mathbf{v}_{d}\\
\mathbf{v}_{1}^{\ast}\\
\vdots\\
\mathbf{v}_{d}^{\ast}%
\end{array}
\right)
\end{equation}
The disadvantage in choosing this route is that one must use a general complex
matrix diagonalization routine instead of complex self adjoint ones, which of
course incurs more computer costs. A further numerical problem now becomes
apparent, although one can prove that if $\mathbf{v}$ is an eigenvector of
$\mathbb{B}$ with eigenvalue $-ie^{i\theta}$ then $\mathbf{v}^{\ast}$ is an
eigenvector with eigenvalue $ie^{i\theta}$ (see Appendix \ref{pairing}), the
eigenvectors, even in absence of degeneracies, are only unique up to an
overall phase factor and thus they do not have to be produced in such complex
conjugate pairs in standard general complex matrix diagonalization routines.
Furthermore if there are degeneracies the eigenvectors corresponding to the
degenerate eigenvalue do not have to be orthogonal, though it is still true
that for normal matrices eigenvectors corresponding to different eigenvalues
are, like the hermitian case, orthogonal. The way to overcome this non
orthogonality and pairing problem is to first L\"{o}wdin symmetric
orthogonalize \cite{Lowdin2} the degenerate eigenvectors corresponding to the
eigenvalue that has a negative real part and then \emph{define }the
eigenvectors associated with the negative of this eigenvalue as the complex
conjugates of these eigenvectors.

Once a paired unitary matrix $\mathbb{V}$ has been produced one can form a
real orthogonal $2d\times2d$ matrix $\mathbb{O}$ by defining its columns by
\begin{align}
\mathbf{o}_{j} &  =\frac{1}{\sqrt{2}}\left(  \mathbf{v}_{j}+\mathbf{v}%
_{j}^{\ast}\right)  =\sqrt{2}\operatorname{Re}\mathbf{v}_{j}\nonumber\\
\mathbf{o}_{j+d} &  =\frac{-i}{\sqrt{2}}\left(  \mathbf{v}_{j}-\mathbf{v}%
_{j}^{\ast}\right)  =\sqrt{2}\operatorname{Im}\mathbf{v}_{j}%
\end{align}
which brings the complex antisymmetric unitary matrix $\mathbb{B}$ to a
canonical form, allowing one to express $\mathbb{B}$ as
\begin{equation}
\mathbb{B}=\left(  \mathbf{o}_{1}\cdots\mathbf{o}_{d}\mathbf{o}_{d+1}%
\mathbf{\cdots o}_{2d}\right)  \left(
\begin{array}
[c]{cccccc}%
0 & 0 & 0 & ie^{i\theta_{1}} & \cdots & 0\\
0 & \ddots & 0 & 0 & \ddots & 0\\
0 & \cdots & 0 & 0 & \cdots & ie^{i\theta_{d}}\\
-ie^{i\theta_{d+1}} & \cdots & 0 & 0 & \cdots & 0\\
0 & \ddots & 0 & 0 & \ddots & 0\\
0 & \cdots & -ie^{i\theta_{2d}} & 0 & \cdots & 0
\end{array}
\right)  \left(
\begin{array}
[c]{c}%
\mathbf{o}_{1}\\
\vdots\\
\mathbf{o}_{d}\\
\mathbf{o}_{d+1}\\
\vdots\\
\mathbf{o}_{2d}%
\end{array}
\right)  \label{reex1}%
\end{equation}
This allows us to construct a unitary matrix $\widetilde{\mathbb{O}}$ that
brings $\mathbb{B}$ to a unique canonical form by defining
\begin{equation}
\left(  \mathbf{\tilde{o}}_{1}\cdots\mathbf{\tilde{o}}_{d}\mathbf{\tilde{o}%
}_{d+1}\mathbf{\cdots\tilde{o}}_{2d}\right)  =\left(  \mathbf{o}_{1}%
\cdots\mathbf{o}_{d}\mathbf{o}_{d+1}\mathbf{\cdots o}_{2d}\right)  \left(
\begin{array}
[c]{cccccc}%
e^{\frac{i\theta_{1}}{2}} & 0 & 0 & 0 & \cdots & 0\\
0 & \ddots & 0 & 0 & \cdots & 0\\
0 & \cdots & e^{\frac{i\theta_{d}}{2}} & 0 & \cdots & 0\\
0 & \cdots & 0 & e^{\frac{i\theta_{d+1}}{2}} & \cdots & 0\\
0 & \cdots & 0 & 0 & \ddots & 0\\
0 & \cdots & 0 & 0 & \cdots & e^{\frac{i\theta_{2d}}{2}}%
\end{array}
\right)
\end{equation}
which leads to
\begin{equation}
\mathbb{B}=\left(  \mathbf{\tilde{o}}_{1}\cdots\mathbf{\tilde{o}}%
_{d}\mathbf{\tilde{o}}_{d+1}\mathbf{\cdots\tilde{o}}_{2d}\right)  \left(
\begin{array}
[c]{cccccc}%
0 & 0 & 0 & 1 & \cdots & 0\\
0 & \ddots & 0 & 0 & \ddots & 0\\
0 & \cdots & 0 & 0 & \cdots & 1\\
-1 & \cdots & 0 & 0 & \cdots & 0\\
0 & \ddots & 0 & 0 & \ddots & 0\\
0 & \cdots & -1 & 0 & \cdots & 0
\end{array}
\right)  \left(
\begin{array}
[c]{c}%
\mathbf{\tilde{o}}_{1}\\
\vdots\\
\mathbf{\tilde{o}}_{d}\\
\mathbf{\tilde{o}}_{d+1}\\
\vdots\\
\mathbf{\tilde{o}}_{2d}%
\end{array}
\right)
\end{equation}
and hence%
\begin{equation}
\mathbf{c}_{\kappa}=\left(  \mathbf{\tilde{o}}_{1}\cdots\mathbf{\tilde{o}%
}_{d_{\kappa}}\mathbf{\tilde{o}}_{d_{\kappa}+1}\mathbf{\cdots\tilde{o}%
}_{2d_{\kappa}}\right)  \left(
\begin{array}
[c]{cc}%
\mathbf{0} & \eta_{\kappa}^{\frac{1}{2}}\mathbb{I}_{d_{\kappa}}\\
-\eta_{\kappa}^{\frac{1}{2}}\mathbb{I}_{d_{\kappa}} & \mathbf{0}%
\end{array}
\right)  \left(
\begin{array}
[c]{c}%
\mathbf{\tilde{o}}_{1}\\
\vdots\\
\mathbf{\tilde{o}}_{d_{\kappa}}\\
\mathbf{\tilde{o}}_{d_{\kappa}+1}\\
\vdots\\
\mathbf{\tilde{o}}_{2d_{\kappa}}%
\end{array}
\right)
\end{equation}
and the full blocked diagonal $\mathbf{c}$ matrix
\begin{align}
\mathbb{Y}^{\dagger}\mathbf{c}\mathbb{Y}^{\ast} &  =\left(
\begin{array}
[c]{ccc}%
\mathbf{c}_{1} & \mathbf{0} & \mathbf{0}\\
\mathbf{0} & \ddots & \mathbf{0}\\
\mathbf{0} & \mathbf{0} & \mathbf{c}_{p}%
\end{array}
\right)  \nonumber\\
&  =\left(
\begin{array}
[c]{ccc}%
\widetilde{\mathbb{O}}_{1} & \mathbf{0} & \mathbf{0}\\
\mathbf{0} & \ddots & \mathbf{0}\\
\mathbf{0} & \mathbf{0} & \widetilde{\mathbb{O}}_{p}%
\end{array}
\right)  \left(
\begin{array}
[c]{ccc}%
\eta_{1}^{\frac{1}{2}}\mathbb{J}_{1} & \mathbf{0} & \mathbf{0}\\
\mathbf{0} & \ddots & \mathbf{0}\\
\mathbf{0} & \mathbf{0} & \eta_{p}^{\frac{1}{2}}\mathbb{J}_{p}%
\end{array}
\right)  \left(
\begin{array}
[c]{ccc}%
\widetilde{\mathbb{O}}_{1} & \mathbf{0} & \mathbf{0}\\
\mathbf{0} & \ddots & \mathbf{0}\\
\mathbf{0} & \mathbf{0} & \widetilde{\mathbb{O}}_{p}%
\end{array}
\right)  ^{t}%
\end{align}
where
\begin{equation}
\mathbb{J}_{\kappa}=\left(
\begin{array}
[c]{cc}%
\mathbf{0} & \mathbb{I}_{d_{\kappa}}\\
-\mathbb{I}_{d_{\kappa}} & \mathbf{0}%
\end{array}
\right)
\end{equation}
This leads to
\begin{equation}
\mathbf{c}=\mathbb{Y}\widetilde{\mathbb{O}}\left(
\begin{array}
[c]{ccc}%
\eta_{1}^{\frac{1}{2}}\mathbb{J}_{1} & \mathbf{0} & \mathbf{0}\\
\mathbf{0} & \ddots & \mathbf{0}\\
\mathbf{0} & \mathbf{0} & \eta_{p}^{\frac{1}{2}}\mathbb{J}_{p}%
\end{array}
\right)  \widetilde{\mathbb{O}}^{t}\mathbb{Y}^{\ast}%
\end{equation}
and hence the unitary matrix $\mathbb{W}^{\dagger}$ whose columns are the
NGSOs by
\begin{equation}
\mathbb{W}^{\dagger}=\mathbb{Y}\widetilde{\mathbb{O}}%
\end{equation}


\section{Discussion\label{discussion}}

We have shown in the preceding that one can form unique canonical coefficients
for a geminal but one must be very careful in the numerical implementation of
the constructions contained in Zuminos' Theorem, i.e. one must constrain the
computed quantities to correspond to the desired mathematical ones. These
unique coefficients have physical significance as they determine the
eigenvalues i.e. occupation numbers of the FORDO of the $2N$-particle
AGP\ state \emph{and} the full Configurational Interaction (CI)\ expansion. A
rather remarkable consequence of this is that the energy the of $2N$- particle
AGP\ state can solely be expressed in terms of the one particle canonical
orbitals and the occupation numbers of the geminal.

\section{Acknowledgements}

We would like thank the support of the Petroleum Research Foundation through a
grant PRF 34859-AC6.

\appendix


\section{Tensor Products\label{tensor}}

In order to use familiar terms we will use the coordinate representation of
one particle states $\left\{  \left\vert \varphi\right\rangle \right\}  $ in
terms of wavefunctions $\varphi\left(  \mathbf{r}\right)  $ to describe the
various tensor products. The coordinate representation is based on the
generalized eigenvectors $\left\{  \left\vert \mathbf{r}\right\rangle
\right\}  $ of the position operator that produce the resolution of the
identity
\begin{equation}
I=\int\left\vert \mathbf{r}\right\rangle \left\langle \mathbf{r}\right\vert
d\mathbf{r}%
\end{equation}
by%
\begin{equation}
\left\vert \varphi\right\rangle =\int\left\vert \mathbf{r}\right\rangle
\left\langle \mathbf{r}|\varphi\right\rangle d\mathbf{r=}\int\left\vert
\mathbf{r}\right\rangle \varphi\left(  \mathbf{r}\right)  d\mathbf{r}%
\end{equation}
The tensor product $\varphi_{1}\otimes\varphi_{2}\left(  \mathbf{r}%
_{1},\mathbf{r}_{2}\right)  $ of two one particle state wavefunctions
$\left\{  \varphi_{1}\left(  \mathbf{r}\right)  ,\varphi_{2}\left(
\mathbf{r}\right)  \right\}  $ is the pointwise product
\begin{equation}
\varphi_{1}\otimes\varphi_{2}\left(  \mathbf{r}_{1},\mathbf{r}_{2}\right)
=\varphi_{1}\left(  \mathbf{r}_{1}\right)  \otimes\varphi_{2}\left(
\mathbf{r}_{2}\right)  =\varphi_{1}\left(  \mathbf{r}_{1}\right)  \varphi
_{2}\left(  \mathbf{r}_{2}\right)
\end{equation}
this product has no particle symmetry but is the easiest to perform
computations such as contractions in. If $\varphi_{1}$ and $\varphi_{2}$ are
normalized than so is $\varphi_{1}\otimes\varphi_{2}$. The two particle state
wavefunction that is the antisymmetric tensor product $\varphi_{1}%
\wedge\varphi_{2}\left(  \mathbf{r}_{1},\mathbf{r}_{2}\right)  $ of two one
particle state wavefunctions $\left\{  \varphi_{1}\left(  \mathbf{r}\right)
,\varphi_{2}\left(  \mathbf{r}\right)  \right\}  $ is given by
\begin{equation}
\varphi_{1}\wedge\varphi_{2}\left(  \mathbf{r}_{1},\mathbf{r}_{2}\right)
=\frac{1}{\sqrt{2}}\left\{  \varphi_{1}\left(  \mathbf{r}_{1}\right)
\otimes\varphi_{2}\left(  \mathbf{r}_{2}\right)  -\varphi_{2}\left(
\mathbf{r}_{1}\right)  \otimes\varphi_{1}\left(  \mathbf{r}_{2}\right)
\right\}  \label{aten}%
\end{equation}
where the factor $\frac{1}{\sqrt{2}}$ has been introduced to maintain
normalization and the RHS of eq. $\left(  \ref{aten}\right)  $ is recognized
as a normalized Slater determinant of the two one particle wavefunctions
$\left\{  \varphi_{1}\left(  \mathbf{r}\right)  ,\varphi_{2}\left(
\mathbf{r}\right)  \right\}  $ i.e.%
\begin{equation}
\varphi_{1}\wedge\varphi_{2}\left(  \mathbf{r}_{1},\mathbf{r}_{2}\right)
=\frac{1}{\sqrt{2}}\left\vert
\begin{array}
[c]{cc}%
\varphi_{1}\left(  \mathbf{r}_{1}\right)  & \varphi_{1}\left(  \mathbf{r}%
_{2}\right) \\
\varphi_{2}\left(  \mathbf{r}_{1}\right)  & \varphi_{2}\left(  \mathbf{r}%
_{2}\right)
\end{array}
\right\vert
\end{equation}


\section{Canonical Form of a Real Antisymmetric Matrix\label{realanti}}

The commuting \emph{real }antisymmetric\emph{ }matrices\emph{ }$\mathbb{N}$
and $\mathbb{M}$ define hermitian matrices
\begin{align}
\mathbb{H}_{N}  &  =i\mathbb{N}\nonumber\\
\mathbb{H}_{M}  &  =i\mathbb{M}%
\end{align}
which can be seen from
\begin{align}
\mathbb{H}_{N}^{\dagger}  &  =\left(  i\mathbb{N}\right)  ^{t\ast}=\left(
i\mathbb{N}^{t}\right)  ^{\ast}=\left(  -i\mathbb{N}\right)  ^{\ast
}=i\mathbb{N=H}_{N}\nonumber\\
\mathbb{H}_{M}^{\dagger}  &  =\left(  i\mathbb{M}\right)  ^{t\ast}=\left(
i\mathbb{M}^{t}\right)  ^{\ast}=\left(  -i\mathbb{M}\right)  ^{\ast
}=i\mathbb{M=H}_{M}%
\end{align}
The hermitian matrices $\mathbb{H}_{N}$ and $\mathbb{H}_{M}$ commute and thus
can be simultaneously diagonalized into real diagonal forms by the same
unitary transformation and hence we show in the following that $\mathbb{N}$
and $\mathbb{M}$ can simultaneously be brought to a canonical antisymmetric
form by using the same real orthogonal transformation that is constructed
according to the following:

Consider the real antisymmetric\emph{ }matrix $\mathbb{A}=\mathbb{A}^{\ast}$
it can be converted into a hermitian purely imaginary matrix, $\mathbb{H}$,
with the properties
\begin{align}
\mathbb{H}  &  =i\mathbb{A}\nonumber\\
\mathbb{H}^{\dagger}  &  =-i\mathbb{A}^{t}=i\mathbb{A}=\mathbb{H}\nonumber\\
\mathbb{H}^{\ast}  &  =-i\mathbb{A}=i\mathbb{A}^{t}\,=\mathbb{H}^{t}%
\end{align}
and
\begin{equation}
\mathbb{H}=i\mathbb{A}=\sum_{1\leq j\leq s}\alpha_{j}\left\{  \left\vert
\mathbf{v}_{j}\right\rangle \left\langle \mathbf{v}_{j}\right\vert -\left\vert
\mathbf{v}_{j}^{\ast}\right\rangle \left\langle \mathbf{v}_{j}^{\ast
}\right\vert \right\}  ;\qquad\alpha_{j}\in\mathbb{R} \label{sa}%
\end{equation}
where $\left\{  \left\vert \mathbf{v}_{j}\right\rangle ,\left\vert
\mathbf{v}_{j}^{\ast}\right\rangle ;1\leq j\leq r\right\}  $ is a complete
orthonormal basis. Eq. $\left(  \ref{sa}\right)  $ is seen by observing that
\begin{equation}
\mathbb{H}\left\vert \mathbf{v}\right\rangle =\alpha\left\vert \mathbf{v}%
\right\rangle \Rightarrow\mathbb{H}^{\ast}\left\vert \mathbf{v}^{\ast
}\right\rangle =\alpha\left\vert \mathbf{v}^{\ast}\right\rangle \Rightarrow
\mathbb{H}^{t}\left\vert \mathbf{v}^{\ast}\right\rangle =\alpha\left\vert
\mathbf{v}^{\ast}\right\rangle \Rightarrow\mathbb{H}\left\vert \mathbf{v}%
^{\ast}\right\rangle =-\alpha\left\vert \mathbf{v}^{\ast}\right\rangle
\end{equation}
It is important to note that the orthonormal basis eigenvectors of
$\mathbb{H}$ are only defined up to an overall phase factors $e^{i\varphi_{j}%
}$ which \emph{do not} preserve the complex conjugate relationship between
eigenvectors associated with paired eigenvalues of opposite sign, however one
can always find complex conjugate eigenvector pairs. These different bases are
related to each other by
\begin{equation}
\left(  \left\vert \mathbf{v}_{1}^{\prime}\right\rangle \cdots\left\vert
\mathbf{v}_{2s}^{\prime}\right\rangle \right)  =\left(  \left\vert
\mathbf{v}_{1}\right\rangle \cdots\left\vert \mathbf{v}_{2s}\right\rangle
\right)  \left(
\begin{array}
[c]{ccc}%
e^{i\varphi_{1}} & \mathbf{0} & \mathbf{0}\\
\mathbf{0} & \ddots & \mathbf{0}\\
\mathbf{0} & \mathbf{0} & e^{i\varphi_{2s}}%
\end{array}
\right)
\end{equation}
Using a basis with complex conjugate symmetry we obtain
\begin{equation}
\mathbb{A}=-\sum_{1\leq j\leq s}i\alpha_{j}\left\{  \left\vert \mathbf{v}%
_{j}\right\rangle \left\langle \mathbf{v}_{j}\right\vert -\left\vert
\mathbf{v}_{j}^{\ast}\right\rangle \left\langle \mathbf{v}_{j}^{\ast
}\right\vert \right\}  ;\qquad\alpha_{i}\in\mathbb{R}%
\end{equation}
We can then define a \emph{real} conb $\left\{  \left\vert \mathbf{s}%
_{j}\right\rangle ;1\leq j\leq\infty\right\}  $ i.e. $\left\vert
\mathbf{s}_{j}\right\rangle =\left\vert \mathbf{s}_{j}^{\ast}\right\rangle $
as
\begin{align}
\left\vert \mathbf{s}_{j}\right\rangle  &  =\frac{1}{\sqrt{2}}\left(
\left\vert \mathbf{v}_{j}\right\rangle +\left\vert \mathbf{v}_{j}^{\ast
}\right\rangle \right)  =\sqrt{2}\operatorname{Re}\left\vert \mathbf{v}%
_{j}\right\rangle \nonumber\\
\left\vert \mathbf{s}_{j+\mathbf{s}}\right\rangle  &  =\frac{-i}{\sqrt{2}%
}\left(  \left(  \left\vert \mathbf{v}_{j}\right\rangle -\left\vert
\mathbf{v}_{j}^{\ast}\right\rangle \right)  \right)  =\sqrt{2}%
\operatorname{Im}\left\vert \mathbf{v}_{j}\right\rangle
\end{align}
so that
\begin{align}
\mathbb{A}\left\{  \frac{1}{\sqrt{2}}\left(  \left\vert \mathbf{v}%
_{j}\right\rangle +\left\vert \mathbf{v}_{j}^{\ast}\right\rangle \right)
\right\}   &  =-i\alpha_{j}\frac{1}{\sqrt{2}}\left(  \left\vert \mathbf{v}%
_{j}\right\rangle -\left\vert \mathbf{v}_{j}^{\ast}\right\rangle \right)
\nonumber\\
&  =\alpha_{j}\left\{  \frac{-i}{\sqrt{2}}\left(  \left(  \left\vert
\mathbf{v}_{j}\right\rangle -\left\vert \mathbf{v}_{j}^{\ast}\right\rangle
\right)  \right)  \right\}
\end{align}
and
\begin{align}
\mathbb{A}\left\{  \frac{-i}{\sqrt{2}}\left(  \left(  \left\vert
\mathbf{v}_{j}\right\rangle -\left\vert \mathbf{v}_{j}^{\ast}\right\rangle
\right)  \right)  \right\}   &  =\frac{-\alpha_{j}}{\sqrt{2}}\left\vert
\mathbf{v}_{j}\right\rangle -\frac{\alpha_{j}}{\sqrt{2}}\left\vert
\mathbf{v}_{j}^{\ast}\right\rangle \nonumber\\
&  =-\alpha_{j}\left\{  \frac{1}{\sqrt{2}}\left(  \left\vert \mathbf{v}%
_{j}\right\rangle +\left\vert \mathbf{v}_{j}^{\ast}\right\rangle \right)
\right\}
\end{align}
thus
\begin{equation}
\mathbb{A=}\sum_{1\leq j\leq\mathbf{s}}\alpha_{j}\left\{  \left\vert
\mathbf{s}_{j}\right\rangle \left\langle \mathbf{s}_{j+\mathbf{s}}\right\vert
-\left\vert \mathbf{s}_{_{j+\mathbf{s}}}\right\rangle \left\langle
\mathbf{s}_{j}\right\vert \right\}
\end{equation}
which displays the desired result. Using the bases $\left\{  \left\vert
\mathbf{s}_{j}\right\rangle ;1\leq j\leq\infty\right\}  $ we can produce the
canonical matrix representation of $\mathbb{A}$
\begin{equation}
\mathbb{A}\left(  \left\vert \mathbf{s}\right\rangle \right)  =\left\vert
\mathbf{s}\right\rangle \left(
\begin{array}
[c]{cc}%
\mathbb{O} & \mathbf{\alpha}\\
-\mathbf{\alpha} & \mathbb{O}%
\end{array}
\right)
\end{equation}
where
\begin{equation}
\mathbf{\alpha=}\left(
\begin{array}
[c]{ccc}%
\alpha_{1} & \mathbb{O} & \mathbb{O}\\
\mathbb{O} & \ddots & \mathbb{O}\\
\mathbb{O} & \mathbb{O} & \alpha_{\mathbf{s}}%
\end{array}
\right)
\end{equation}
and we have the expansion
\begin{equation}
\mathbb{A}=\left\vert \mathbf{s}\right\rangle \left(
\begin{array}
[c]{cc}%
\mathbb{O} & \mathbf{\alpha}\\
-\mathbf{\alpha} & \mathbb{O}%
\end{array}
\right)  \left\langle \mathbf{s}\right\vert \label{reex}%
\end{equation}
The row vector of basis vectors, $\left(  \left\vert \mathbf{s}_{1}%
\right\rangle \cdots\left\vert \mathbf{s}_{\mathbf{s}}\right\rangle \left\vert
\mathbf{s}_{\mathbf{s}+1}\right\rangle \cdots\left\vert \mathbf{s}%
_{\mathbf{s}+\mathbf{s}}\right\rangle \right)  $, can be viewed as a real
orthogonal matrix whose columns are the real basis vectors $\left\{
\left\vert \mathbf{s}_{1}\right\rangle ,\cdots\left\vert \mathbf{s}%
_{\mathbf{s}}\right\rangle ,\left\vert \mathbf{s}_{\mathbf{s}+1}\right\rangle
,\cdots\left\vert \mathbf{s}_{\mathbf{s}+\mathbf{s}}\right\rangle \right\}  $.

The complex basis $\left\{  \left\vert \mathbf{v}_{j}\right\rangle ,\left\vert
\mathbf{v}_{j}^{\ast}\right\rangle ;1\leq j\leq r\right\}  $ diagonalizes
$\mathbb{A}$ to give
\begin{equation}
\mathbb{A}\left(  \left\vert \mathbf{v}\right\rangle ,\left\vert
\mathbf{v}^{\ast}\right\rangle \right)  =\left(  \left\vert \mathbf{v}%
\right\rangle ,\left\vert \mathbf{v}^{\ast}\right\rangle \right)  \left(
\begin{array}
[c]{cc}%
-i\mathbf{\alpha} & \mathbb{O}\\
\mathbb{O} & i\mathbf{\alpha}%
\end{array}
\right)
\end{equation}
and produces the expansion%
\begin{equation}
A=\left(  \left\vert \mathbf{v}\right\rangle ,\left\vert \mathbf{v}^{\ast
}\right\rangle \right)  \left(
\begin{array}
[c]{cc}%
-i\mathbf{\alpha} & \mathbb{O}\\
\mathbb{O} & i\mathbf{\alpha}%
\end{array}
\right)  \left(
\begin{array}
[c]{c}%
\left\langle \mathbf{v}\right\vert \\
\left\langle \mathbf{v}^{\ast}\right\vert
\end{array}
\right)  \label{ccex}%
\end{equation}
The row vector of basis vectors, $\left(  \left\vert \mathbf{v}_{1}%
\right\rangle \cdots\left\vert \mathbf{v}_{\mathbf{s}}\right\rangle \left\vert
\mathbf{v}_{1}^{\ast}\right\rangle \cdots\left\vert \mathbf{v}_{\mathbf{s}%
}^{\ast}\right\rangle \right)  $, can be viewed as an unitary matrix whose
columns are the basis vectors $\left\{  \left\vert \mathbf{v}_{1}\right\rangle
\cdots\left\vert \mathbf{v}_{\mathbf{s}}\right\rangle \left\vert
\mathbf{v}_{1}^{\ast}\right\rangle \cdots\left\vert \mathbf{v}_{\mathbf{s}%
}^{\ast}\right\rangle \right\}  $. Note that the expansions of eqs. $\left(
\ref{reex}\right)  $ and $\left(  \ref{ccex}\right)  $ both give $\mathbb{A}$
thus in eq. $\left(  \ref{reex}\right)  $ all the imaginary terms cancel.

\section{Complex Conjugate Pairing\label{pairing}}

Let $\mathbb{A}$ be a complex antisymmetric matrix that is normal then
\begin{equation}
\left[  \mathbb{A},\mathbb{A}^{\dagger}\right]  =0
\end{equation}
and $\mathbb{A}$ and $\mathbb{A}^{\dagger}$ are simultaneously diagonalizable
and the spectral resolution gives
\begin{subequations}
\begin{align}
\mathbb{A}  &  =\sum_{1\leq j\leq r}\alpha_{j}\left\vert \mathbf{v}%
_{j}\right\rangle \left\langle \mathbf{v}_{j}\right\vert \\
\mathbb{A}^{\dagger}  &  =\sum_{1\leq j\leq r}\alpha_{j}^{\ast}\left\vert
\mathbf{v}_{j}\right\rangle \left\langle \mathbf{v}_{j}\right\vert
\label{adag}%
\end{align}
so there exist simultaneous eigenvectors $\left\{  \mathbf{v}_{j}\right\}  $
to $\mathbb{A}$ and $\mathbb{A}^{\dagger}$ corresponding to complex conjugate
eigenvalues $\left\{  \alpha,\alpha^{\ast}\right\}  $ i.e.%
\end{subequations}
\begin{subequations}
\begin{align}
\mathbb{A}\left\vert \mathbf{v}_{j}\right\rangle  &  =\alpha\left\vert
\mathbf{v}_{j}\right\rangle \nonumber\\
\mathbb{A}^{\dagger}\left\vert \mathbf{v}_{j}\right\rangle  &  =\alpha
_{j}^{\ast}\left\vert \mathbf{v}_{j}\right\rangle
\end{align}
as $\mathbb{A}$ is antisymmetric we obtain from eq. $\left(  \ref{adag}%
\right)  $ that
\end{subequations}
\begin{equation}
\mathbb{A}^{\ast}\left\vert \mathbf{v}_{j}\right\rangle =-\alpha_{j}^{\ast
}\left\vert \mathbf{v}_{j}\right\rangle
\end{equation}
which gives by complex conjugation that%
\begin{equation}
\mathbb{A}\left\vert \mathbf{v}_{j}^{\ast}\right\rangle =-\alpha_{j}\left\vert
\mathbf{v}_{j}^{\ast}\right\rangle
\end{equation}
\bibliographystyle{apsrev}
\bibliography{agpopt,bw,novel}



\end{document}